\section{Results}\label{sec:results}

\subsection{Experiment 1: replication (RQ1)}
\begin{table*}[!tbp]
  \centering
  \caption{Edge-IIoTset, test set: data as published vs.\ with the placeholder spelling unified (``--'': not run).}
  \label{tab:exp1}
  \gentable{exp1_replication}
\end{table*}
On the data as published, PPFLE + small BERT reaches
\res{exp1.bert_small.e1_ppfle.edge_iiotset_raw.multiclass.accuracy} accuracy, against the paper's
\res{exp1.paper.accuracy}. The number therefore reproduces (RQ1: \res{hyp.RQ1.verdict}).
% CHECK hyp.RQ1.verdict
As \cref{sec:leak} shows, however, it only reproduces because of the leak. On the fixed data the same model reaches
\res{exp1.bert_small.e1_ppfle.edge_iiotset.multiclass.accuracy} accuracy and
\res{exp1.bert_small.e1_ppfle.edge_iiotset.multiclass.macro_f1} macro-F1. A plain XGBoost does better, with
\res{exp1.xgboost.none.edge_iiotset.multiclass.accuracy} accuracy and
\res{exp1.xgboost.none.edge_iiotset.multiclass.macro_f1} macro-F1.

\subsection{Experiment 2: within-dataset benchmark (RQ2, H1, H2)}
\begin{table*}[!tbp]
  \centering
  \caption{Test macro-F1 (multiclass), mean $\pm$ std over \res{seeds.distilbert} seeds.}
  \label{tab:exp2}
  \gentable{exp2_within_macro_f1}
\end{table*}
\Cref{tab:exp2} shows the within-dataset results (as a chart: \cref{fig:exp2}, \cref{app:more}). XGBoost is as good as or better than every text model on
every dataset (H2: \res{hyp.H2.verdict}, \res{hyp.H2.held} of \res{hyp.H2.total} datasets). On NF-UNSW-NB15 it leads
the best text model by \resdiff{exp2.xgboost.none.nf_unsw_nb15.macro_f1}{exp2.distilbert.e3_kv.nf_unsw_nb15.macro_f1}
macro-F1, and on NF-BoT-IoT by
\resdiff{exp2.xgboost.none.nf_bot_iot.macro_f1}{exp2.distilbert.e3_kv.nf_bot_iot.macro_f1}, just above our tie
margin. On NF-ToN-IoT the best
text model comes within our tie margin. On NF-CSE-CIC-IDS2018, DistilBERT with E3 is even slightly ahead
(\resdiff{exp2.distilbert.e3_kv.nf_cse_cic_ids2018.macro_f1}{exp2.xgboost.none.nf_cse_cic_ids2018.macro_f1}), but
that is still a tie.

On Edge-IIoTset (fixed data), XGBoost reaches \res{exp2.xgboost.none.edge_iiotset.macro_f1}. The best small BERT
(\res{best.edge_iiotset.bert_small.encoding}) reaches \res{best.edge_iiotset.bert_small.value}, and DistilBERT with
\res{best.edge_iiotset.distilbert.encoding} reaches \res{best.edge_iiotset.distilbert.value}. That is the highest
value in the column, but its lead over XGBoost
(\resdiff{best.edge_iiotset.distilbert.value}{exp2.xgboost.none.edge_iiotset.macro_f1}) is within the tie margin.
% CHECK hyp.H2.verdict and its per-dataset checks in results/analysis_summary.md, best.<d>.method (NF-CSE-CIC-IDS2018
% and Edge-IIoTset: DistilBERT E3), best.nf_unsw_nb15.distilbert.encoding

For DistilBERT, readable key=value text (E3) is the best encoding on every dataset, with the sentence template (E4)
usually close behind. For the small BERT, the best encoding changes from dataset to dataset. Raw PPFLE is never the
best encoding for either model and is usually among the weakest. Its best model is within
\res{cfg.analysis.hypotheses.well_margin} of the best method on only \res{hyp.H1.held} of \res{hyp.H1.total}
datasets (H1: \res{hyp.H1.verdict}). Seed variance also matters. On NF-CSE-CIC-IDS2018, the small BERT with E1 varies
by $\pm$\res{exp2.bert_small.e1_ppfle.nf_cse_cic_ids2018.macro_f1.std} between seeds, which is enough to change the
order of the encodings.
% CHECK hyp.H1.verdict; best.<d>.distilbert.encoding = E3 on all five datasets; E1 never ranks first

\paragraph{Error analysis.} Most of the errors come from a few classes (\cref{app:errors}). On NF-UNSW-NB15, the
\texttt{generic} class is hard for every method. XGBoost reaches a per-class F1 of
\res{perclass.xgboost.none.nf_unsw_nb15.generic} on it, and the small BERT with E1 only
\res{perclass.bert_small.e1_ppfle.nf_unsw_nb15.generic}. The rare \texttt{worms} class is found by XGBoost
(\res{perclass.xgboost.none.nf_unsw_nb15.worms}) and by DistilBERT with E3
(\res{perclass.distilbert.e3_kv.nf_unsw_nb15.worms}), but the small BERT with E1 misses it
(\res{perclass.bert_small.e1_ppfle.nf_unsw_nb15.worms}). On NF-CSE-CIC-IDS2018, every method labels most
\texttt{infiltration} flows as benign (XGBoost:
\res{confusion.xgboost.none.nf_cse_cic_ids2018.infiltration.benign}\,\%). This suggests that these flows cannot be told
apart from benign traffic using the NetFlow features. On Edge-IIoTset, no class stays below a per-class F1 of
\res{cfg.analysis.errors.hard_class_f1} for every method. The strongest methods (XGBoost, DistilBERT with E3 or E4, and the
small BERT with E4) all make the same mistakes: application-layer attacks that look like benign traffic. XGBoost,
for example, labels
\res{confusion.xgboost.none.edge_iiotset.ddos_http.benign}\,\% of \texttt{ddos\_http} flows and
\res{confusion.xgboost.none.edge_iiotset.uploading.benign}\,\% of \texttt{uploading} flows as benign. With binned
PPFLE (E2), both language models instead confuse attack classes with each other, for example \texttt{uploading} with
\texttt{password}, probably because the bins merge values that would separate these attacks.
% CHECK hard.edge_iiotset and the top mix-ups per method (computed from results/run_details)

\subsection{Experiment 3: cross-dataset transfer (RQ3, H3)}
\begin{table*}[!tbp]
  \centering
  \caption{Mean over the \res{exp3.n_datasets}$\times$\res{exp3.n_datasets} grid: same dataset (diagonal) vs.\
    transfer (off-diagonal).}
  \label{tab:exp3}
  \gentable{exp3_cross}
\end{table*}
When a model is trained and tested on the same dataset, every method detects attacks almost perfectly. XGBoost, for
example, reaches a diagonal macro-F1 of \res{exp3.xgboost.none.diag.macro_f1}. Transfer to another dataset is much
harder (\cref{tab:exp3}; all pairs: \cref{fig:exp3}, \cref{app:more}). The average off-diagonal macro-F1 of XGBoost is only
\res{exp3.xgboost.none.off.macro_f1}, and some pairs collapse completely. XGBoost trained on NF-CSE-CIC-IDS2018 and
tested on NF-BoT-IoT reaches just \res{exp3.xgboost.none.nf_cse_cic_ids2018.nf_bot_iot.macro_f1}.

Binning helps both language models. With E2 instead of E1, the small BERT improves from
\res{exp3.bert_small.e1_ppfle.off.macro_f1} to \res{exp3.bert_small.e2_ppfle_bins.off.macro_f1}, and DistilBERT from
\res{exp3.distilbert.e1_ppfle.off.macro_f1} to \res{exp3.distilbert.e2_ppfle_bins.off.macro_f1}. That is the highest
off-diagonal score of any method. Readable text (E3, E4) helps the small BERT but not DistilBERT (H3:
\res{hyp.H3.verdict}). For DistilBERT, both readable encodings even
transfer worse than raw PPFLE (\res{exp3.distilbert.e3_kv.off.macro_f1} and
\res{exp3.distilbert.e4_sentence.off.macro_f1} against \res{exp3.distilbert.e1_ppfle.off.macro_f1}). DistilBERT with E2 also catches more attacks
than XGBoost, with an attack recall of \res{exp3.distilbert.e2_ppfle_bins.off.attack_recall} against
\res{exp3.xgboost.none.off.attack_recall}, at a similar false-positive rate
(\res{exp3.distilbert.e2_ppfle_bins.off.fpr} against \res{exp3.xgboost.none.off.fpr}; \cref{app:more}).
% CHECK hyp.H3.verdict and the E2 ranking in exp3_cross

\subsection{Does pretraining help readable encodings? (H4)}
H4 predicts that DistilBERT beats the small BERT on readable encodings only. The data show gains on readable text,
for example \resdiff{exp2.distilbert.e3_kv.nf_unsw_nb15.macro_f1}{exp2.bert_small.e3_kv.nf_unsw_nb15.macro_f1}
macro-F1 for E3 on NF-UNSW-NB15. But they also show gains on hashes, for example
\resdiff{exp3.distilbert.e2_ppfle_bins.off.macro_f1}{exp3.bert_small.e2_ppfle_bins.off.macro_f1} off-diagonal
macro-F1 for E2 (H4: \res{hyp.H4.verdict}). So pretraining helps beyond reading English, but not everywhere. On
Edge-IIoTset both models score the same on hashes, \res{exp2.distilbert.e1_ppfle.edge_iiotset.macro_f1} against
\res{exp2.bert_small.e1_ppfle.edge_iiotset.macro_f1} for E1 and
\res{exp2.distilbert.e2_ppfle_bins.edge_iiotset.macro_f1} against
\res{exp2.bert_small.e2_ppfle_bins.edge_iiotset.macro_f1} for E2, while DistilBERT does gain on readable text. Model
size is a confounder here (\res{params.distilbert.millions}\,M against \res{params.bert_small.millions}\,M
parameters). A larger model that starts from pretrained weights may simply fit token statistics better.
Across datasets the pattern even reverses. There, DistilBERT gains on hashes
(\res{exp3.distilbert.e1_ppfle.off.macro_f1} against \res{exp3.bert_small.e1_ppfle.off.macro_f1} for E1) but not on
readable text (\res{exp3.distilbert.e3_kv.off.macro_f1} against \res{exp3.bert_small.e3_kv.off.macro_f1} for E3, and
\res{exp3.distilbert.e4_sentence.off.macro_f1} against \res{exp3.bert_small.e4_sentence.off.macro_f1} for E4).
% CHECK hyp.H4.verdict and its cross-dataset checks (E1/E2 gain, E3/E4 no gain)

\subsection{Experiment 4: is PPFLE private? (RQ4)}\label{sec:privacy}
\begin{table*}[!tbp]
  \centering
  \caption{Dictionary attack on the NetFlow test splits: share of tokens and of whole rows recovered (cumulative
    levels).}
  \label{tab:privacy}
  \gentable{privacy_attack}
\end{table*}
PPFLE uses an unsalted hash. Anyone who knows the scheme can therefore hash candidate values and compare them with
the tokens. Our attacker knows the scheme and the public NetFlow feature list (column names and types), but not the
data. The attack runs in three cumulative steps (\cref{tab:privacy}; \cref{fig:privacy}, \cref{app:privacy}).

First, we try every possible value of the \res{privacy.n_low_cardinality} low-cardinality columns, such as ports,
protocols, TCP flags, TTLs and application codes. This alone reverses
\res{privacy.e1_ppfle.enumerate.token_share.min}--\res{privacy.e1_ppfle.enumerate.token_share.max}\,\% of all E1
tokens. Second, we try every whole number up to \res{privacy.bruteforce.max} for the remaining numeric columns. This
raises the share to
\res{privacy.e1_ppfle.bruteforce.token_share.min}--\res{privacy.e1_ppfle.bruteforce.token_share.max}\,\% and already
recovers many whole rows exactly, for example \res{privacy.nf_cse_cic_ids2018.e1_ppfle.bruteforce.row_share}\,\% of
the NF-CSE-CIC-IDS2018 test flows. Third, we add a dictionary of the values seen in the other public datasets. Token
recovery then reaches \res{privacy.e1_ppfle.public.token_share.min}--\res{privacy.e1_ppfle.public.token_share.max}\,\%,
and full rows reach \res{privacy.nf_cse_cic_ids2018.e1_ppfle.public.row_share}\,\% on NF-CSE-CIC-IDS2018 and
\res{privacy.nf_ton_iot.e1_ppfle.public.row_share}\,\% on NF-ToN-IoT. The values that resist are fractional
bytes-per-second rates and very large throughput counters (\cref{fig:privacycols}, \cref{app:privacy}).

E2 is even easier to attack. Every binned column has only \res{cfg.encoders.binning.n_bins} possible tokens, so
enumeration alone reverses \res{privacy.e2_ppfle_bins.enumerate.token_share.min}\,\% of the tokens. The attacker then
learns the bin, which is a quantile range, rather than the exact value. We checked every recovered token against the
true value.

The attack also works on SecurityBERT's own dataset (\cref{tab:privacyedge}, \cref{app:privacy}). There is no independent public dataset
with the same columns as Edge-IIoTset, so we only use the first two steps. The text columns are guessed from public
vocabularies (HTTP methods and versions, MQTT protocol names) plus the two empty-field spellings. On the data as
published, this reverses \res{privacy_edge.edge_iiotset_raw.e1_ppfle.bruteforce.token_share}\,\% of the PPFLE tokens
and \res{privacy_edge.edge_iiotset_raw.e1_ppfle.bruteforce.row_share}\,\% of whole rows. What stays hidden is mostly
the TCP sequence and acknowledgement numbers, whose values are larger than our brute-force range. Of the raw
acknowledgement numbers, for example, only \res{privacy_edge.col.e1_ppfle.bruteforce.tcp.ack_raw}\,\% are recovered.

Deterministic hashing of low-entropy fields is pseudonymisation, not privacy protection. A keyed hash (an HMAC with a
secret key) would block this dictionary attack. It would still reveal which values are equal, though, and the model
would still learn from exact values that rarely recur in another network.
